% ===========================================================================
\section{Pipeline y camino de datos}
\label{sec:pipeline}

\diagrama{datapath.png}{Camino de datos}{Las cinco etapas con sus bloques, los cuatro registros
de segmentación, las señales de control y los caminos de forwarding y de redirección. Es el
diseño final, con el comparador de igualdad para los saltos (\decision{044}). Fuente editable:
\repo{docs/diagramas/datapath.drawio}.}{fig:datapath}

El procesador es un RV32I clásico de cinco etapas, con datos y direcciones de 32~bits
(\repo{docs/spec/pipeline.md}). La figura~\ref{fig:datapath} muestra el camino de datos completo.
El módulo superior es \repo{rtl/pipeline/riscv_core.v}, que instancia una etapa por módulo
(\repo{if_stage}, \repo{id_stage}, \repo{ex_stage}, \repo{mem_stage}, \repo{wb_stage}), los
cuatro registros de segmentación y las unidades de forwarding y de riesgos.

\begin{center}
\small
\begin{tabular}{@{}l p{6.2cm} p{5.6cm}@{}}
  \toprule
  \textbf{Etapa} & \textbf{Qué hace} & \textbf{Bloques} \\
  \midrule
  IF  & Mantiene el PC y lee la instrucción & \texttt{pc\_reg}, sumador $+4$, mux de próximo PC, memoria de programa \\
  ID  & Decodifica, lee registros, arma el inmediato y el control; detecta la carga-uso & banco de registros, generador de inmediatos, unidad de control \\
  EX  & Opera, calcula direcciones y decide los saltos & muxes de forwarding, ALU, sumador $PC+imm$, comparador de igualdad \\
  MEM & Lee o escribe la memoria de datos & memoria de datos con \emph{byte write enable} \\
  WB  & Extiende el dato de un load y escribe \texttt{rd} & extensión, mux de write-back \\
  \bottomrule
\end{tabular}
\end{center}

\subsection{Registros de segmentación}

Cada registro de segmentación guarda todo lo que una instrucción necesita en las etapas
siguientes. Todos tienen las mismas tres entradas de control, con prioridad
\texttt{rst} $>$ \texttt{en} $>$ \texttt{flush}:

\begin{center}
\small
\begin{tabular}{@{}l l@{}}
  \toprule
  \textbf{Entrada} & \textbf{Efecto} \\
  \midrule
  \texttt{rst}   & Carga una burbuja \\
  \texttt{en}    & En 0 conserva el contenido (stall o core detenido por la Debug Unit) \\
  \texttt{flush} & En 1 carga una burbuja en lugar de la instrucción entrante \\
  \bottomrule
\end{tabular}
\end{center}

Una \textbf{burbuja} es un registro con \texttt{valid = 0} y todas las señales de control con
efectos en 0 (escritura de registro y de memoria, salto y halt). Sus campos de datos no se
limpian, porque no importan. Que \texttt{en} gane sobre \texttt{flush} garantiza que un stall o el
core detenido nunca pierden la instrucción retenida.

\begin{center}
\small
\begin{tabular}{@{}l r l@{}}
  \toprule
  \textbf{Registro} & \textbf{Bits} & \textbf{Campos} \\
  \midrule
  IF/ID  &  65 & \texttt{valid}, \texttt{pc}, \texttt{instr} \\
  ID/EX  & 161 & \texttt{valid}, \texttt{pc}, \texttt{rs1\_data}, \texttt{rs2\_data}, \texttt{imm}, \texttt{rs1}, \texttt{rs2}, \texttt{rd}, \texttt{funct3} y 14 bits de control \\
  EX/MEM & 109 & \texttt{valid}, \texttt{pc}, \texttt{result}, \texttt{store\_data}, \texttt{rd}, \texttt{funct3} y 4 bits de control \\
  MEM/WB & 110 & \texttt{valid}, \texttt{pc}, \texttt{result}, \texttt{read\_data}, \texttt{rd}, \texttt{funct3}, \texttt{offset} y 3 bits de control \\
  \midrule
  \textbf{Total} & \textbf{445} & 381 en flip-flops y 64 en los registros de salida de las BRAM \\
  \bottomrule
\end{tabular}
\end{center}

Los campos se llaman igual en la especificación, en el RTL (\texttt{id\_ex\_alu\_ctrl},
\texttt{ex\_mem\_reg\_write}, \ldots), en el modelo de referencia y en el banco de pruebas, que los
lee por jerarquía. El orden de los campos define también cómo se empaquetan en la traza del
modelo y en el volcado que la Debug Unit manda a la PC (decisiones 020 y 023): los tres
coinciden bit a bit.

\subsection{Etapas}

\paragraph{IF.} El próximo PC se elige con esta prioridad:
\begin{lstlisting}[style=verilog,numbers=none]
pc_next = redirect   ? redirect_pc :   // salto tomado en EX
          stop_fetch ? pc_reg      :   // HALT en vuelo (seccion 3.1)
                       pc_reg + 4;
\end{lstlisting}
La memoria de programa es una BRAM de lectura sincrónica: la dirección entra en un flanco y la
instrucción sale en el siguiente, ya registrada. Por eso el registro de salida de la BRAM
\emph{es} el campo \texttt{instr} de IF/ID, y su \emph{enable} es el mismo \texttt{en\_if\_id}
(\decision{009}).

\paragraph{ID.} El banco de registros tiene 32 registros en flip-flops, dos puertos de lectura
asíncronos y uno de escritura manejado por WB (\decision{007}). Si WB escribe el registro que ID
está leyendo en el mismo ciclo, la lectura entrega el dato nuevo (\emph{bypass} interno): resuelve
la dependencia a distancia~3 sin forwarding. Un tercer puerto de lectura, sin bypass, es para el
volcado de la Debug Unit. La unidad de control decodifica por opcode y decide también la operación
de la ALU, así que a EX solo viajan 4~bits de \texttt{alu\_ctrl} (\decision{026}).

\paragraph{EX.} Los operandos pasan por los muxes de forwarding (sección~\ref{sec:forwarding}).
Los saltos se resuelven acá (\decision{012}), con un sumador propio para $PC + imm$ y un
comparador de igualdad que no pasa por la ALU (\decision{044}):
\begin{lstlisting}[style=verilog,numbers=none]
eq          = (rs1_val == rs2_val);              // operandos con forwarding
taken       = branch & (eq ^ funct3[0]);         // beq: eq = 1, bne: eq = 0
redirect    = taken | jal | jalr;
redirect_pc = jalr ? {alu_out[31:1], 1'b0} : pc + imm;
\end{lstlisting}
Alcanza con igualdad porque el subset solo tiene \texttt{beq} y \texttt{bne}. El comparador sacó
del camino crítico el mux del operando~B, el desplazador y el mux de resultado de la ALU
(sección~\ref{sec:timing}).

\paragraph{MEM.} La memoria de datos es otra BRAM con \emph{byte write enable}: \texttt{sb} y
\texttt{sh} escriben solo los bytes que corresponden, sin leer y reescribir la palabra
(\decision{016}). El dato leído sale registrado al final del ciclo y es el campo
\texttt{read\_data} de MEM/WB.

\paragraph{WB.} Extiende el dato del load según \texttt{funct3} y el desplazamiento dentro de la
palabra (\texttt{lb}, \texttt{lh}, \texttt{lbu}, \texttt{lhu}), elige entre el resultado de EX y
el dato de memoria y escribe \texttt{rd}. También actualiza \texttt{halted} cuando llega un HALT
válido.

\subsection{Convenciones del RTL}

Todo el RTL sigue las convenciones de la \decision{010}: cabecera con función, puertos, latencia y
valor de reset; prefijos \texttt{i\_}/\texttt{o\_}, \texttt{\_reg}/\texttt{\_next} y \texttt{u\_};
reset sincrónico y activo en alto solo en control, \texttt{valid} y PC; un solo dominio de reloj;
codificaciones solo en \repo{rtl/common/riscv_defs.vh}. Ningún diseño infiere latches
(\texttt{Register as Latch = 0} en la utilización).
